\chapter{MARCO REFERENCIAL}

\section{Marco legal y regulaciones}

En el Perú, las operaciones de minería subterránea se rigen por el Reglamento de Seguridad y Salud Ocupacional en Minería, aprobado mediante el Decreto Supremo N$^\circ$ 024-2016-EM y sus modificatorias del Ministerio de Energía y Minas (MINEM), cuya fiscalización compete a SUNAFIL y OSINERGMIN~\cite{MINEM_DS024_2016, OSINERGMIN_BoletinGSM_2020}. Este reglamento obliga al titular minero a verificar condiciones atmosféricas respirables antes y durante el ingreso del personal a cualquier frente subterráneo. En particular, el artículo 246 exige monitorear y mantener por debajo de los Límites de Exposición Ocupacional del Anexo 15 las concentraciones de monóxido de carbono ($\text{CO}$), dióxido de carbono ($\text{CO}_2$), gases nitrosos ($\text{NO}_2$) y demás contaminantes, asegurando además un mínimo volumétrico de oxígeno ($\text{O}_2$)~\cite{MINEM_DS024_2016, OSINERGMIN_DS024_Compilado}. De forma complementaria, los artículos 30, 250 y 254 ordenan paralizar las labores o clausurar accesos ante deficiencias de ventilación o acumulaciones de humo diésel y gases de voladura. Estas disposiciones sustentan legalmente el desarrollo de un sistema robótico cooperativo que ejecute la primera inspección ambiental bajo protocolos de no ingreso (\textit{no-entry}), evitando exponer a supervisores y cuadrillas a frentes ciegos no verificados~\cite{MINEM_DS024_2016, OSINERGMIN_BoletinGSM_Ventilacion_2020}.

\section{Identificación de normas y estándares de ingeniería}

Las especificaciones metrológicas y constructivas del sistema toman como referencia estándares técnicos internacionales aplicables a instrumentación ambiental y robótica móvil. Para el monitoreo de gases tóxicos y asfixiantes, la Directiva Europea 2017/164 sobre Valores Límite de Exposición Profesional Indicativos (IOELV) y los criterios internacionales de armonización de límites ocupacionales establecen umbrales de referencia para $\text{CO}$, $\text{CO}_2$, $\text{NO}$ y $\text{NO}_2$ que orientan la resolución, el rango dinámico y el tiempo de respuesta ($t_{90}$) requeridos en los transductores~\cite{EU_ExposureDirective_2020, Deveau2015_OEL}. Con el fin de evitar lecturas sesgadas por polvo fino o condensación en galerías húmedas, la selección e integración de sensores químicos debe contemplar filtrado físico en la admisión de muestra, acondicionamiento analógico y rutinas periódicas de calibración con gas patrón~\cite{Cepa2023_sensors, EU_ExposureDirective_2020}.

En cuanto a la protección electromecánica y la arquitectura de software, el diseño considera estándares de estanqueidad e interoperabilidad modular. Las carcasas proyectadas para resguardar la electrónica de control y distribución eléctrica toman como meta constructiva el grado IP65 de la norma IEC 60529 frente al ingreso de polvo abrasivo y goteo de agua en socavón~\cite{AdvancesMining_2023, OSINERGMIN_BoletinGSM_2020}. Por su parte, la navegación del hexápodo primario se estructura sobre el marco ROS 2 empleando el estándar de comunicaciones DDS (\textit{Data Distribution Service}) y árboles estandarizados de transformaciones espaciales (\texttt{tf2}), mientras que la intercomunicación inalámbrica con los módulos embebidos del hexápodo secundario se resuelve mediante tramas UDP con marcas temporales consistentes para preservar la trazabilidad en entornos sin señal GNSS~\cite{siciliano2010robotics, KimChoi_2025_MineMapping}.

\section{Identificación de restricciones generales}

El dimensionamiento físico, energético y económico de ambos hexápodos está acotado por las condiciones geométricas de las labores subterráneas y por los recursos disponibles para el prototipado experimental~\cite{MINEM_DS024_2016, OSINERGMIN_BoletinGSM_2020}. En el aspecto geométrico, las galerías de avance, cruceros y rampas presentan pendientes de hasta $12\%$ y secciones reducidas donde el D.S. N$^\circ$ 024-2016-EM exige preservar pasos peatonales libres de al menos $1.0\text{ m}$~\cite{MINEM_DS024_2016}. Ello obliga a contener la envergadura en reposo por debajo de $0.50\text{ m}$ para no obstruir vías de tránsito ni quedar bloqueados entre hastiales. Asimismo, en el hexápodo secundario accionado por micromotores DC N20, cada extremidad de tres grados de libertad se sujeta a una proporción cinemática $1x : 2x : 3x$ (coxa de $3.0\text{ cm}$, fémur de $6.0\text{ cm}$ y tibia de $9.0\text{ cm}$) que reduce la coxa al mínimo constructivo admisible por el alojamiento del motor, acotando el brazo de palanca horizontal y el par flector en la cadera~\cite{InsectInspired_2020, Wu2024_SpiderLegDesign}.

A estas fronteras mecánicas se suman limitaciones de propagación electromagnética y de presupuesto de manufactura. La atenuación de radiofrecuencia provocada por la roca encajonante impide depender de un enlace continuo hacia superficie, exigiendo procesamiento embarcado y registro local no volátil. En el plano financiero, la construcción íntegra de las dos plataformas, tanto la unidad primaria con servomotores MG996R como la secundaria con motores N20, se ajusta a un presupuesto acotado de investigación universitaria. Esta restricción orienta la fabricación estructural hacia manufactura aditiva FDM en ácido poliláctico (PLA) con $8\%$ de relleno en rejilla (\textit{grid}) y tres perímetros de contorno, configuración adoptada tras comprobar experimentalmente en banco que las piezas impresas con solo dos paredes sufrían fisuras prematuras bajo los esfuerzos transmitidos por los ejes motrices~\cite{Anderson2017_PLA, MahovicPoljacek2025_PLA}.

\section{Compromiso ético y responsabilidad profesional con la seguridad y el bienestar social}

La concepción de este sistema cooperativo responde al deber de emplear la ingeniería mecatrónica para proteger la integridad física de los trabajadores en actividades de alto riesgo, en concordancia con el principio de prevención de la Ley N$^\circ$ 29783 (Ley de Seguridad y Salud en el Trabajo), el D.S. N$^\circ$ 024-2016-EM y el Objetivo de Desarrollo Sostenible 3 (Salud y Bienestar)~\cite{MINEM_DS024_2016}. Sustituir el ingreso presencial de inspectores a labores ciegas por robots de verificación previa ataca directamente una de las causas recurrentes de mortalidad por inhalación de gases en la minería nacional~\cite{MINEM_Yanaquihua_2023}. Esta responsabilidad profesional exige también rigor y honestidad metrológica: reportar con fidelidad los alcances y límites reales del prototipo experimental, garantizar registros digitales íntegros sin alteración de lecturas y priorizar arquitecturas de fabricación accesible que permitan, en etapas futuras de maduración industrial, acercar herramientas de monitoreo remoto a operaciones mineras de pequeña y mediana escala~\cite{OSINERGMIN_BoletinGSM_2020}.
